\bcpchapters{Private Baptism of Children}{The Ministration of Private Baptism of Children}

\begin{rubric}
The Curates of every parish shall often admonish the people, that they defer not the Baptism of their children longer than the first or second Sunday after their birth, or other holy-day falling between, unless a great and reasonable cause, to be approved by the Curate.
\end{rubric}

\begin{rubric}
And also they shall warn the people that without great cause and necessity they procure not their children to be baptized at home in their houses.
But when need shall compel them so to do, then Baptism shall be administered on this fashion:
\end{rubric}

\begin{rubric}
First, let the Minister of the Parish (or, in his absence, any other lawful Minister that can be procured) with them that are present, call upon God, and say the Lord's Prayer, and so many of the Collects appointed to be said before in the form of Publick Baptism, as the time and present exigence will suffer.
And then, the Child being named by some one that is present, the Minister shall pour Water upon it, saying these words;
\end{rubric}

\lett{N}{ I} baptize thee In the Name of the Father, and of the Son, and of the Holy Ghost.
\inlineresp{Amen.}

\begin{rubric}
Then all kneeling down, the Minister shall give thanks unto God, and say,
\end{rubric}

\lett{W}{e} yield thee hearty thanks, most merciful Father, that it hath pleased thee to regenerate this infant with thy Holy Spirit, to receive him for thine own Child by adoption, and to incorporate him into thy holy Church.
And we humbly beseech thee to grant, that as he is now made partaker of the death of thy Son, so he may be also of his resurrection; and that finally, with the residue of thy Saints, he may inherit thine everlasting kingdom; through the same thy Son Jesus Christ our Lord.
\inlineresp{Amen.}

\newpage

\begin{rubric}
And let them not doubt, but that the Child so baptized is lawfully and sufficiently baptized, and ought not to be baptized again.
Yet nevertheless, if the Child, which is after this sort baptized, do afterward live, it is expedient that it be brought into the Church, to the intent that, if the Minister of the same Parish did himself baptize that Child, the Congregation may be certified of the true form of Baptism, by him privately before used: In which case he shall say thus,
\end{rubric}

\lett{I}{ certify} you, that according to the due and prescribed Order of the Church, \changeable{at such a time}, and \changeable{at such a place}, before divers witnesses I baptized this child.

\begin{rubric}
But if the Child were baptized by any other lawful Minister, then the Minister of the Parish, where the Child was born or christened, shall examine and try whether the Child be lawfully baptized, or no.
In which case, if those that bring any Child to the Church do answer, that the same Child is already baptized, then shall the Minister examine them further, saying,
\end{rubric}

\lett{B}{y} whom was this child baptized?
Who was present when this child was baptized?
Because some things essential to this Sacrament may happen to be omitted through fear or haste, in such times of extremity; therefore I demand further of you,
With what matter was this child baptized?
With what words was this child baptized?

\begin{rubric}
And if the Minister shall find by the answers of such as bring the Child, that all things were done as they ought to be; then shall not he christen the Child again, but shall receive him as one of the flock of true Christian people, saying thus,
\end{rubric}

\lett{I}{ certify} you, that in this case all is well done, and according unto due order, concerning the baptizing of this child; who being born in original sin, and in the wrath of God, is now, by the laver of Regeneration in Baptism, received into the number of the children of God, and heirs of everlasting life: for our Lord Jesus Christ doth not deny his grace and mercy unto such Infants, but most lovingly doth call them unto him, as the holy Gospel doth witness to our comfort on this wise.

\newpage

\begin{leftbar}

\begin{center}
Hear the words of the Gospel, written by Saint Mark, in the tenth chapter, beginning at the thirteenth verse.
\end{center}

\inrub{Answer.} Glory be to thee, O Lord.

\end{leftbar}

\lett{T}{hey} brought young children to Christ, that he should touch them; and his disciples rebuked those that brought them.
But when Jesus saw it, he was much displeased, and said unto them, Suffer the little children to come unto me, and forbid them not; for of such is the kingdom of God.
Verily I say unto you, Whosoever shall not receive the kingdom of God as a little child, he shall not enter therein.
And he took them up in his arms, put his hands upon them, and blessed them.

\begin{leftbar}
\inrub{Answer.} Praise be to thee, O Christ.
\end{leftbar}

\begin{rubric}
After the Gospel is read, the Minister shall make this brief exhortation upon the words of the Gospel.
\end{rubric}

\lett{B}{eloved}, ye hear in this Gospel the words of our Saviour Christ, that he commanded the children to be brought unto him; how he blamed those that would have kept them from him; how he exhorted all men to follow their innocency.
Ye perceive how by his outward gesture and deed he declared his good will toward them; for he embraced them in his arms, he laid his hands upon them, and blessed them.
Doubt ye not therefore, but earnestly believe, that he hath likewise favourably received this present Infant; that he hath embraced him with the arms of his mercy: and (as he hath promised in his holy Word) will give unto him the blessing of eternal life, and make him partaker of his everlasting kingdom.
Wherefore, we being thus persuaded of the good will of our heavenly Father, declared by his Son Jesus Christ, towards this infant, let us faithfully and devoutly give thanks unto him, and say the Prayer which the Lord himself taught us:

\lett{O}{ur} Father, which art in heaven, Hallowed be thy Name; Thy kingdom come; Thy will be done; In earth as it is in heaven. Give us this day our daily bread. And forgive us our trespasses, As we forgive them that trespass against us. And lead us not into temptation; But deliver us from evil.

\inlineresp{Amen.}

\newpage

\lett{A}{lmighty} and everlasting God, heavenly Father, we give thee humble thanks, that thou hast vouchsafed to call us to the knowledge of thy grace, and faith in thee: Increase this knowledge, and confirm this faith in us evermore.
Give thy Holy Spirit to this infant, that he, being born again, and being made an heir of everlasting salvation, through our Lord Jesus Christ, may continue thy servant, and attain thy promise; through the same our Lord Jesus Christ thy Son, who liveth and reigneth with thee and the Holy Spirit, now and for ever.
\inlineresp{Amen.}

\begin{rubric}
Then shall the Priest demand the Name of the Child; which being by the Godfathers and Godmothers pronounced, the Minister shall say,
\end{rubric}

\lett{D}{ost} thou, in the name of this child, renounce the devil and all his works, the vain pomp and glory of this world, with all covetous desires of the same, and the carnal desires of the flesh, so that thou wilt not follow, nor be led by them?

\inrub{Answer.} I renounce them all.

\lett{D}{ost} thou believe in God the Father Almighty, Maker of heaven and earth?
And in Jesus Christ his only-begotten Son our Lord?
And that he was conceived by the Holy Ghost; born of the Virgin Mary; that he suffered under Pontius Pilate, was crucified, dead, and buried; that he went down into hell, and also did rise again the third day; that he ascended into heaven, and sitteth at the right hand of God the Father Almighty; and from thence shall come again at the end of the world, to judge the quick and the dead?
And dost thou believe in the Holy Ghost; the holy Catholick Church; the Communion of Saints; the Remission of sins; the Resurrection of the flesh; and everlasting life after death?

\inrub{Answer.} All this I stedfastly believe.

\lett{W}{ilt} thou then obediently keep God's holy will and commandments, and walk in the same all the days of thy life?

\inrub{Answer.} I will.

\newpage

\begin{rubric}
Then shall the Priest say,
\end{rubric}

\lett{W}{e} receive this child into the congregation of Christ's flock, *and do sign \changeable{him} with the sign of the Cross, in token that hereafter \changeable{he} shall not be ashamed to confess the faith of Christ crucified, and manfully to fight under his banner, against sin, the world, and the devil; and to continue Christ's faithful soldier and servant unto \changeable{his} life's end.
\begin{secrubric}
* Here the Priest shall make a Cross upon the child's forehead.
\end{secrubric}

\begin{rubric}
Then shall the Priest say,
\end{rubric}

\lett{S}{eeing} now, dearly beloved brethren, that this child is by Baptism regenerate, and grafted into the body of Christ's Church, let us give thanks unto Almighty God for these benefits; and with one accord make our prayers unto him, that \changeable{he} may lead the rest of \changeable{his} life according to this beginning.

\begin{rubric}
Then shall the Priest say,
\end{rubric}

\lett{W}{e} yield thee most hearty thanks, most merciful Father, that it hath pleased thee to regenerate \changeable{this infant} with thy Holy Spirit, to receive \changeable{him} for thine own Child by adoption, and to incorporate \changeable{him} into thy holy Church.
And humbly we beseech thee to grant, that \changeable{he} being dead unto sin, and living unto righteousness, and being buried with Christ in his death, may crucify the old man, and utterly abolish the whole body of sin; and that, as \changeable{he} is made partaker of the death of thy Son, \changeable{he} may also be partaker of his resurrection; so that finally, with the residue of thy holy Church, \changeable{he} may be \changeable{an inheritor} of thine everlasting kingdom; through Jesus Christ our Lord.
\inlineresp{Amen.}

\begin{rubric}
Then, all standing up, the Minister shall make this exhortation to the Godfathers and Godmothers.
\end{rubric}

\lett{F}{orasmuch} as \changeable{this child hath} promised by you \changeable{his} sureties to renounce the devil and all his works, to believe in God, and to serve him: ye must remember, that it is your parts and duties to see that \changeable{this infant} be taught, so soon as \changeable{he} shall be able to learn, what a solemn vow, promise, and profession \changeable{he hath} made by you.
And that \changeable{he} may know these things the better, ye shall call upon \changeable{him} to hear Sermons; and chiefly ye shall provide, that \changeable{he} may learn the Creed, the Lord's Prayer, and the Ten Commandments, in the vulgar tongue, and all other things which a Christian ought to know and believe to \changeable{his} soul's health; and that this child may be virtuously brought up to lead a godly and a Christian life; remembering alway, that Baptism doth represent unto us our profession; which is, to follow the example of our Saviour Christ, and to be made like unto him; that, as he died, and rose again for us, so should we, who are baptized, die from sin, and rise again unto righteousness; continually mortifying all our evil and corrupt affections, and daily proceeding in all virtue and godliness of living.

\begin{rubric}
But if they which bring the infant to the Church do make such uncertain answers to the Priest's questions, as that it cannot appear that the child was baptized with water, In the Name of the Father, and of the Son, and of the Holy Ghost, (which are essential parts of Baptism,) then let the Priest baptize it in the form before appointed for Publick Baptism of Infants; saving that at the dipping of the child in the Font, or the pouring of water upon it, he shall use this form of words.
\end{rubric}

\lett{I}{f} thou art not already baptized, \changeable{N.} I baptize thee In the Name of the Father, and of the Son, and of the Holy Ghost.
\inlineresp{Amen.}
